\section*{Capítulo \ref{ch:b1:organism-environment} --- O organismo: um sistema em interação com o ambiente}

\begin{solution}{exo:b1:organism-environment:1}
Matéria (um mamífero come alimento e expira dióxido de carbono; uma
planta absorve dióxido de carbono e libera oxigênio), energia (um
mamífero capta a energia química do alimento e libera calor; uma planta
capta luz e libera calor), informação (um mamífero enxerga um predador;
uma planta se curva na direção da luz).
\end{solution}

\begin{solution}{exo:b1:organism-environment:2}
$S/V = 3/r$ cai pela metade. A captação através da superfície cresce
como $r^2$ e o consumo como $r^3$; a partir de certo raio a superfície
não consegue mais alimentar o volume, e por isso as células ficam na
faixa do micrômetro ou dobram a superfície.
\end{solution}

\begin{solution}{exo:b1:organism-environment:3}
Manutenção de uma variável do \omterm{def:b1:organism-environment:milieu}{meio interno} perto de um \omterm{def:b1:organism-environment:homeostasis}{valor de
referência} por \omterm{def:b1:organism-environment:homeostasis}{retroalimentação negativa}. Sensor: os termorreceptores da
pele e do encéfalo; centro integrador: o hipotálamo, que guarda o \omterm{def:b1:organism-environment:homeostasis}{valor
de referência} de \qty{37}{\celsius}; efetores: os músculos que tremem,
os vasos da pele e as glândulas sudoríparas.
\end{solution}

\begin{solution}{exo:b1:organism-environment:4}
Lagarto: de $1.0$ a $0.25$, um fator $4$ (duas duplicações para
\qty{20}{\celsius}). Camundongo: de $1.0$ a $3.4$, um fator de cerca de
$3.4$.
\end{solution}

\begin{solution}{exo:b1:organism-environment:5}
Aresta de \qty{1}{mm}: precisa de $\qty{1}{nmol/s}$
($V = \qty{1}{mm^3}$); oferta $2\times 6 = \qty{12}{nmol/s}$: sim.
Aresta de \qty{10}{mm}: precisa de \qty{1000}{nmol/s}, oferta
$2\times 600 = \qty{1200}{nmol/s}$: por pouco. Limite:
$2\times 6L^2 = L^3$ dá $L = \qty{12}{mm}$.
\end{solution}

\begin{solution}{exo:b1:organism-environment:6}
Cavalo: $3.4\times 500^{0.75} = \qty{360}{W}$, \qty{0.72}{W/kg}.
Musaranho: $3.4\times 0.005^{0.75} = \qty{0.064}{W}$, \qty{13}{W/kg}.
Por dia: $3\times 0.064\times 86400 = \qty{16.6}{kJ}$, ou seja,
\qty{3.3}{g} de alimento, dois terços da própria massa.
\end{solution}

\begin{solution}{exo:b1:organism-environment:7}
$\qty{80}{W} = \qty{4.8}{kJ/min}$, ou seja, \qty{0.24}{L} de
$\mathrm{O_2}$ por minuto; por dia $80\times 86400 = \qty{6.9}{MJ}$.
\qty{2000}{kcal} = \qty{8.36}{MJ} por dia = \qty{97}{W}.
\end{solution}

\begin{solution}{exo:b1:organism-environment:8}
Camundongo: $100\times(0.02/70)^{0.75} = \qty{0.22}{m^2}$; elefante:
$100\times(4000/70)^{0.75} = \qty{2080}{m^2}$. Esferas: camundongo
$r = \qty{1.7}{cm}$, $S = \qty{0.0036}{m^2}$ (o pulmão é sessenta vezes
a superfície do corpo); elefante $r = \qty{0.98}{m}$,
$S = \qty{12}{m^2}$ (o pulmão é cento e setenta vezes). Quanto maior o
animal, mais o pulmão dele precisa ser dobrado em relação ao exterior.
\end{solution}

\begin{solution}{exo:b1:organism-environment:9}
Os efetores só agem depois que o sensor registrou um desvio, de modo que
é preciso haver um desvio antes que ele seja corrigido; a correção então
ultrapassa ligeiramente o alvo antes de o sensor informá-lo. A largura
da faixa é fixada pela sensibilidade do sensor (o menor desvio
detectável), pelo atraso entre a detecção e o efeito, e pela intensidade
dos efetores.
\end{solution}

\begin{solution}{exo:b1:organism-environment:10}
$V = \qty{2e-5}{m^3}$, $r = \qty{1.68}{cm}$, $S = 4\pi r^2 =
\qty{3.6e-3}{m^2}$. Perda: $5\times 3.6\times 10^{-3}\times 37 =
\qty{0.66}{W}$. Produção basal: \qty{0.18}{W}, quatro vezes menos. O
camundongo precisa aumentar a produção (tremor, alimentação: três a
quatro vezes a basal), reduzir $h$ (ninho, aglomeração) ou deixar a
temperatura cair (torpor).
\end{solution}

\begin{solution}{exo:b1:organism-environment:11}
Perda $\propto M^{2/3}$, produção $\propto M^{3/4}$: a razão
produção/perda $\propto M^{1/12}$ cresce com a massa, de modo que os
animais pequenos têm a menor produção por unidade de superfície e os
grandes, a maior. Esfera de massa $M$:
$S = 4\pi(3M/4000\pi)^{2/3} = \num{0.0483}\,M^{2/3}$ (SI); perda
$= 5\times 20\times 0.0483\,M^{2/3} = 4.83\,M^{2/3}$; a igualdade com
$3.4\,M^{3/4}$ dá $M^{1/12} = 1.42$, $M \approx \qty{67}{kg}$. Abaixo
disso, o animal em repouso passa frio sem isolamento; acima, ele se
aquece e precisa dissipar calor.
\end{solution}

\begin{solution}{exo:b1:organism-environment:12}
Não é mais simples: é organizada de outro modo. As trocas da planta são
externas (superfícies de folha e de raiz espalhadas pelo meio) e as
células dela ficam expostas a um meio que ela não regula; a
``regulação'' dela é o crescimento e a forma, que ajustam as próprias
superfícies. O mamífero interioriza as trocas por trás de um líquido
regulado, de modo que as células dele nunca veem o exterior. A
\omterm{def:b1:organism-environment:homeostasis}{homeostase} é uma propriedade do \omterm{def:b1:organism-environment:milieu}{meio interno} do animal, e não uma medida
de complexidade: as células da planta regulam o próprio conteúdo, os
estômatos dela regulam a perda de água, e o corpo dela é uma \omterm{prop:b1:organism-environment:surfaces}{superfície
de troca} em ampliação contínua --- uma estratégia, não uma ausência.
\end{solution}

\begin{solution}{pb:b1:organism-environment:1}
\textbf{1.} $V = M/\rho$: camundongo \qty{2e-5}{m^3} (\qty{20}{cm^3}),
elefante \qty{4}{m^3}.
\textbf{2.} $r = (3V/4\pi)^{1/3}$: camundongo \qty{1.68}{cm}, elefante
\qty{0.985}{m}.
\textbf{3.} $S = 4\pi r^2$: camundongo \qty{3.55e-3}{m^2}
(\qty{35.5}{cm^2}), elefante \qty{12.2}{m^2}.
\textbf{4.} $S/V = 3/r$: camundongo \qty{178}{m^{-1}}, elefante
\qty{3.05}{m^{-1}}; razão $58$.
\textbf{5.} $S/V \propto 1/r \propto M^{-1/3}$, de modo que a razão é
$(M_e/M_m)^{1/3} = (2\times 10^5)^{1/3} = 58.5$.
\textbf{6.} Maior para os dois (a esfera tem a menor superfície para o
seu volume); o camundongo, com a cauda, as orelhas e os membros finos,
se afasta mais.
\textbf{7.} $P = 10\times S\times 20$: camundongo \qty{0.71}{W},
elefante \qty{2440}{W}.
\textbf{8.} Camundongo $3.4\times 0.02^{0.75} = \qty{0.18}{W}$;
elefante $3.4\times 4000^{0.75} = \qty{1710}{W}$.
\textbf{9.} Perda de \qty{2440}{W} contra produção de \qty{1710}{W} em
repouso: a pele nua é aproximadamente adequada, e o elefante corre mais
risco de superaquecer que de esfriar quando está ativo ou num ar mais
quente.
\textbf{10.} Perda de \qty{0.71}{W} contra \qty{0.18}{W}: quatro vezes
demais. A pelagem precisa baixar $h$ para cerca de
\qty{2.5}{W.m^{-2}.K^{-1}}.
\textbf{11.} O $S/V$ dele é o menor de qualquer mamífero terrestre; as
orelhas acrescentam um terço à superfície e, sendo finas e irrigadas,
são radiadores que despejam calor quando o corpo produz mais do que
\qty{12}{m^2} conseguem perder.
\textbf{12.} $S/V \propto M^{-1/3}$ cresce sem limite à medida que $M$
cai, enquanto $B/M \propto M^{-1/4}$ cresce mais devagar; abaixo de
alguns gramas nenhuma pelagem e nenhuma taxa de alimentação cobrem a
perda. Deixar a temperatura cair à noite (torpor) elimina o termo
$T_b - T_a$ nas horas sem alimento.
\textbf{13.} $2B\times\qty{86400}{s}$: camundongo \qty{31}{kJ},
elefante \qty{295}{MJ}.
\textbf{14.} Camundongo $31/18 = \qty{1.7}{g}$; elefante
$295\,000/5 = \qty{59}{kg}$.
\textbf{15.} Camundongo $1.7/20 = 8.6\%$ da massa dele; elefante
$59/4000 =
1.5\%$: o camundongo come seis vezes mais por grama.
\textbf{16.} Razão entre os comprimentos do intestino $35/0.4 = 88$;
razão entre os raios $0.985/0.0168
= 59$. O intestino cresce mais rápido que o tamanho linear do corpo
(aqui, aproximadamente como $M^{0.36}$), o que aumenta a superfície
absortiva do elefante em relação ao volume dele.
\textbf{17.} $f_0 = 600\times 0.02^{1/4} = \qty{226}{min^{-1}}$;
$\tau_0 = 2/0.02^{1/4} = \qty{5.3}{yr}$.
\textbf{18.} $f = 226\times 4000^{-1/4} = \qty{28}{min^{-1}}$; $\tau =
5.3\times 4000^{1/4} = \qty{42}{yr}$.
\textbf{19.} Camundongo $600\times 60\times 24\times 365\times 2 =
\num{6.3e8}$; elefante $28\times 60\times 24\times 365\times 42 =
\num{6.2e8}$.
\textbf{20.} Batimentos por vida $= f\tau \propto M^{-1/4}M^{1/4} = M^0$:
independentes da massa, perto de um bilhão.
\textbf{21.} $f = 226\times 70^{-1/4} = \qty{78}{min^{-1}}$ (perto de
70); $\tau = 5.3\times 70^{1/4} = \qty{15}{yr}$, muito abaixo de 80: a
longevidade humana é o ponto fora da curva, de três a cinco vezes a
tendência dos mamíferos, por razões (encéfalo, sociabilidade, medicina)
alheias à lei de escala.
\textbf{22.} $160\times 4000^{-1/4}/0.02^{-1/4} = 160\times 0.0473 =
\qty{7.6}{min^{-1}}$; batimentos por respiração $600/160 = 28/7.6 = 3.75$
para os dois, já que as duas frequências compartilham o expoente.
\textbf{23.} $B/M = 3.4\,M^{3/4}/M = 3.4\,M^{-1/4}$; expoente $-1/4$.
\textbf{24.} Camundongo $0.18/0.02 = \qty{9.0}{W/kg}$; elefante
$1710/4000 =
\qty{0.43}{W/kg}$.
\textbf{25.} Razão $9.0/0.43 = 21$, igual a $(M_e/M_m)^{1/4} =
(2\times 10^5)^{1/4} = 21.1$: cada grama do camundongo queima vinte e
uma vezes mais rápido que cada grama do elefante --- a razão entre as
\omterm{def:b1:organism-environment:allometry}{taxas metabólicas} específicas por massa é a raiz quarta da razão entre
as massas.
\end{solution}
